\documentclass[14pt,a4paper]{extarticle}
\usepackage[T2A]{fontenc}
\usepackage[utf8]{inputenc}
\usepackage[english,russian]{babel}
\usepackage{lmodern}
\usepackage{geometry}
\usepackage{array}
\usepackage{tabularx}
\usepackage{longtable}
\usepackage{ltablex}
\usepackage{booktabs}
\usepackage{hyperref}
\usepackage{indentfirst}
\usepackage{microtype}
\usepackage{graphicx}
\usepackage{setspace}
\usepackage{listings}
\usepackage{graphicx}
\usepackage{caption}
\usepackage{float}
\usepackage{tocloft}
\usepackage{titlesec}
\usepackage{newtxtext,newtxmath}

\titleformat{\section}        % форматируем \section
  {\normalfont\bfseries}      % стиль: обычный + жирный
  {\thesection.}              % номер секции с точкой
  {2mm}                       % расстояние между номером и текстом
  {}                           % перед заголовком можно добавить что-то
\titlespacing*{\section}{1.5cm}{*3}{*1}  % {отступ слева}{отступ сверху}{отступ снизу}

\titleformat{\subsection}        % форматируем \section
  {\normalfont\bfseries}      % стиль: обычный + жирный
  {\thesubsection.}              % номер секции с точкой
  {2mm}                       % расстояние между номером и текстом
  {}                           % перед заголовком можно добавить что-то
\titlespacing*{\subsection}{1.5cm}{*3}{*1}  % {отступ слева}{отступ сверху}{отступ снизу}

\captionsetup{justification=centering}
\renewcommand{\ttdefault}{cmtt}
\keepXColumns
\setlength{\parindent}{3em}
\newgeometry{left=2.5cm,right=2.5cm,top=3cm,bottom=2.5cm}

\renewcommand{\contentsname}{\hspace*{4em}{\smallСодержание}}
\setlength{\cftbeforesecskip}{0.5em}
\renewcommand{\cftsecleader}{\cftdotfill{\cftdotsep}}
\renewcommand{\cftsecaftersnum}{.\hspace{2mm}}

\title{Отчёт по "Управление ИТ-проектами"}
\author{Выполнил студент группы М-РИТ-25-1, Станиславчук Сергей}

\begin{document}

\setstretch{1.5}
\noindent
\begin{minipage}[t]{0.01\textwidth}
    \includegraphics{~/programmer/++Programmer/5_1/LSTU.jpg}
\end{minipage}%
\begin{minipage}[c]{1.1\linewidth}
    \centering
    {\small
    МИНИСТЕРСТВО НАУКИ И ВЫСШЕГО ОБРАЗОВАНИЯ РФ \\[1em]
    ФЕДЕРАЛЬНОЕ ГОСУДАРСТВЕННОЕ БЮДЖЕТНОЕ \\
    ОБРАЗОВАТЕЛЬНОЕ УЧРЕЖДЕНИЕ ВЫСШЕГО ОБРАЗОВАНИЯ \\
    «ЛИПЕЦКИЙ ГОСУДАРСТВЕННЫЙ ТЕХНИЧЕСКИЙ УНИВЕРСИТЕТ»
    }
\end{minipage}%

\vspace{3em}

\noindent
\hspace{1em}%
\begin{minipage}[t]{0.8\textwidth}
\textbf{Институт} \hspace{2em} компьютерных наук \\
\textbf{Кафедра} \hspace{2em} автоматизированных систем управления
\end{minipage}


\vspace{6em}

\begin{center}
\textbf{Практическая работа №6} \\
по дисциплине \\ <<Управление ИТ-проектами>>
\end{center}

\vspace{6em}

\noindent
\begin{tabular*}{\textwidth}{@{\extracolsep{\fill}} l c r }
    Студент М-РИТ-25-1 & \makebox[4cm]{\hrulefill} & Станиславчук С.М. \\[2em]
Руководитель & \\
К.т.н., доцент & \makebox[4cm]{\hrulefill} & Ведищев В.В. \\
\end{tabular*}

\vfill
\begin{center}
    Липецк, 2026
\end{center}

\thispagestyle{empty}

\newpage
\tableofcontents
\newpage

\section{Планирование версий программного продукта}

\subsection{Фазы процесса разработки}

Процесс разработки программного продукта разбивается на последовательные фазы, каждая из которых соответствует определённому состоянию системы. На начальном этапе осуществляется формирование требований и проектирование архитектуры. Далее следует фаза активной разработки, в рамках которой реализуется основной функционал системы.

После реализации ключевых компонентов начинается этап тестирования, направленный на выявление ошибок и проверку соответствия системы требованиям. Завершающей фазой является выпуск стабильной версии продукта и его последующа1я эксплуатация с возможностью внесения доработок и исправлений.

Каждая фаза сопровождается выпуском промежуточных версий, что позволяет контролировать прогресс разработки и своевременно выявлять проблемы.

\subsection{Планирование и построение дерева версий}

Для управления развитием системы используется схема версионирования вида X.Y.Z, где каждая компонента отражает уровень изменений в продукте.

\textbf{Мажорная версия (X.Y.Z)} изменяется в случае значительных изменений
системы. К таким изменениям относятся переработка архитектуры, внедрение новых
ключевых функций или изменение существующих интерфейсов, несовместимое с
предыдущими версиями. Выпуск мажорной версии означает качественный переход
системы на новый уровень.

\textbf{Минорная версия (X.Y.Z)} увеличивается при добавлении нового
функционала без нарушения совместимости с предыдущими версиями. Это могут быть
новые возможности системы, расширение существующих функций или улучшения, не
влияющие на основную архитектуру.

\textbf{Патч-версия (X.Y.Z)} увеличивается при внесении незначительных
изменений, таких как исправление ошибок, оптимизация производительности или
незначительные доработки интерфейса. Такие изменения не затрагивают
функциональность системы с точки зрения пользователя.

Дерево версий реализуется через систему ветвления. Основной веткой является \texttt{main}, содержащая стабильную версию системы. Параллельно используется ветка \texttt{develop}, в которой происходит интеграция новых функций перед выпуском релиза.

Функциональность разрабатывается в отдельных ветках, создаваемых от \texttt{develop}. Например:
\begin{itemize}
\item ветки вида \texttt{feature/chat} — реализация текстовых сообщений;
\item \texttt{feature/voice} — реализация голосовой связи (UDP, RTP);
\item \texttt{feature/file-transfer} — передача файлов;
\item \texttt{feature/auth} — аутентификация и управление пользователями;
\item \texttt{feature/ui} — разработка пользовательского интерфейса.
\end{itemize}

Исправления ошибок выполняются в ветках типа \texttt{bugfix/*}, которые также создаются от \texttt{develop} и после проверки сливаются обратно.

Для критических исправлений в уже выпущенных версиях используются ветки \texttt{hotfix/*}, создаваемые от \texttt{main}. После исправления изменения интегрируются как в \texttt{main}, так и в \texttt{develop}.

Перед выпуском версии формируется ветка \texttt{release/X.Y.0}, в которой выполняется финальное тестирование и подготовка релиза. После этого ветка сливается в \texttt{main}, и версия помечается тегом.

\subsection{Средства контроля версий}

В качестве системы контроля версий используется Git, при котором для каждой роли или задачи создаются отдельные ветки. Это позволяет изолировать изменения и упростить процесс их интеграции.

Основная ветка содержит стабильную версию продукта, в то время как разработка нового функционала ведётся в отдельных ветках. После завершения работы изменения проходят проверку и объединяются с основной веткой.

Для обозначения стабильных релизов используются теги, соответствующие номеру версии. Это позволяет фиксировать состояние системы на момент выпуска и упрощает возврат к предыдущим версиям при необходимости.

Ответственность за управление версиями и выпуск релизов возлагается на бэкенд-разработчика, который контролирует целостность системы, корректность версионирования и процесс интеграции изменений.
\end{document}
